\selectlanguage{french}
\Language{fr}{Français}{Guide d'utilisation}
\firsttopic{Ouvrir, mettre en veille, revenir}
Inkline exécute un shell local sur reMarkable 2. Utilisez un Type Folio ou un
clavier USB, travaillez sur la tablette ou connectez-vous par SSH ou Goblin Mosh.
Le texte et les images sont affichés en niveaux de gris.

\begin{notice}
\key{Opt+RightAlt+T} ouvre Inkline directement sur la tablette.\par
\key{Bouton d'alimentation :} appuyez puis relâchez pour mettre en veille ;
appuyez de nouveau pour réveiller. Les shells et les tampons de l'éditeur restent en RAM.
\end{notice}
La veille suspend la tablette sans quitter Inkline ni enregistrer ses réglages
sur la mémoire flash. Le réseau peut nécessiter une reconnexion. L'appui de
réveil ne déclenche pas une nouvelle mise en veille. Après une mise à jour,
quittez puis rouvrez Inkline pour activer la nouvelle gestion de ce bouton.

Pour revenir aux carnets, touchez \key{Quit} ou appuyez sur \key{Ctrl+Shift+Q},
puis confirmez. \code{exit} ferme un terminal ; fermer le dernier rend la main
aux carnets. La barre du bas propose Escape, historique haut/bas, Quit et Settings.
\code{clear} efface le petit guide de démarrage et le logo Goblin.

Si une application de raccourci occupe l'écran, \key{Opt+RightAlt+T} revient à Inkline.
En urgence, maintenez \key{Opt+RightAlt+Backspace pendant deux secondes}.
Cette récupération ferme tous les terminaux et perd leur travail non enregistré.
Après un redémarrage de la tablette, l'interface habituelle des carnets reste celle par défaut.

RightAlt désigne la touche Alt/Opt à droite de la barre d’espace. Maintenez-la avec la touche Opt distincte, puis appuyez sur T, sans Ctrl. Super seul reste disponible pour vos raccourcis.

\ProjectLinks

\topic{Touches et emplacements de terminal}
Le Type Folio possède une touche \key{Opt} distincte entre Ctrl et Alt,
et une touche \key{Alt/Opt droite}. Leurs fonctions sont différentes.

\begin{keytable}
\key{Touches} & \key{Action} \\ \midrule
Opt distincte + 1--0 & F1--F10 \\
Alt/Opt droite + 1--9 & Choisir l'emplacement 1--9 \\
Alt/Opt droite + Left / Right & Emplacement précédent / suivant \\
Alt/Opt droite + Tab & Escape \\
Alt/Opt droite + Up / Down & PageUp / PageDown \\
Alt/Opt droite + Backspace & Confirmer la fermeture de tous les terminaux \\
L'une des touches Alt + Space & Clavier Unicode \\
Alt/Opt droite + C / V & Copier la sélection / coller \\
Ctrl+Shift+B & Masquer / afficher la barre du bas \\
\end{keytable}

Sur le \key{Type Folio américain}, Alt/Opt droite+\key{0} saisit \code{+}.
Alt/Opt droite avec la \key{touche moins juste à gauche de Backspace} saisit
\code{=}. Avec la méthode de saisie désactivée, Shift+6 donne un accent
circonflexe littéral : \code{c^2} reste \code{c^2}. Le zoom n'a pas de raccourci clavier.

Les touches de fonction ordinaires des claviers USB fonctionnent normalement.
Le raccourci Opt distinct utilise Meta si le clavier USB le fournit.

Il existe au plus \key{neuf emplacements indépendants}. Choisir un emplacement
vide ouvre un shell. L'indicateur compte les terminaux ouverts : si seuls 1 et 9
sont ouverts, le 9 affiche \key{Terminal 2 / 2 · slot 9}. Fermer un shell libère
son emplacement sans renuméroter les raccourcis.

Chaque terminal a son shell, son répertoire, sa composition en cours, ses images
et 500 lignes d'historique. Les réglages d'affichage et de saisie sont communs.

\topic{Réglages et affichage}
Touchez \key{Settings}, le cinquième bouton du bas. Si la barre est masquée,
faites \key{Alt+Space}, puis \key{F2}, ou touchez Settings dans le clavier Unicode.
Un choix rempli indique l'état actif ; un contour pointillé indique le focus.
Tab déplace le focus ; les flèches et Entrée modifient le choix.

\key{Caps Lock :} choisissez Control (par défaut) ou Caps Lock.
\key{Taille du texte :} pincez avec deux doigts ou utilisez les boutons moins/plus.
Un mouvement lent règle la taille par pas d'un pixel, de \key{6 à 48 pixels}.

\key{Text darkness :} 50\% correspond au rendu original. Une valeur supérieure
assombrit les contours. \key{Minimum contrast} sépare la luminance du texte et
du fond lorsque leurs couleurs sont proches ; sa valeur par défaut est 35\%.

\begin{choices}
\key{Mode e-paper} & \key{Compromis} \\ \midrule
Fast & Attente minimale ; forme d'onde d'animation ; davantage de rémanence \\
Balanced & Forme d'onde d'interface ; regroupement sur 32 ms \\
Crisp (par défaut) & Gris plus nets ; forme d'onde de contenu ; plus lent \\
Mono & Noir et blanc rapide, sans gris dans les images \\
Saver & Regroupement sur 120 ms ; moins de rafraîchissements en rafale \\
\end{choices}

Seules les lignes modifiées sont redessinées, mais la dalle a son propre délai.
L'économie de batterie de Saver n'a pas été mesurée. Une mise à jour conserve
un mode explicitement enregistré.

\begin{notice}
Tous les réglages restent en \key{RAM} et sont enregistrés ensemble à la sortie
normale d'Inkline. Lever le doigt, fermer Settings ou mettre en veille n'écrit
pas sur la flash. Sans modification, aucune écriture. Un plantage, un arrêt
forcé ou une coupure peut perdre les changements non enregistrés.
\end{notice}

\topic{Unicode et méthodes de saisie}
\key{Alt+Space} ouvre le clavier Unicode. Touchez un caractère pour l'insérer ;
le clavier reste ouvert. Touchez \key{Done} ou appuyez sur Escape pour fermer.

Les catégories regroupent lettres, marques, nombres, ponctuation, mathématiques/
monnaies, symboles, espaces, caractères de format et usage privé. Faites glisser
verticalement, ou utilisez les flèches, PageUp/PageDown, Home/End et la molette.
\key{Tab} change de catégorie ; Entrée insère le caractère sélectionné.

Saisissez une valeur Unicode hexadécimale, par exemple \code{03bb}, puis Entrée
pour insérer λ. Backspace corrige la valeur. Le cercle pointillé des marques
combinatoires sert d'exemple ; seule la marque est insérée. Espaces et caractères
de format portent des étiquettes. Le catalogue utilise les tables Unicode de Qt,
sans contrôles, valeurs non attribuées, substituts ni non-caractères. Certains
glyphes nécessitent une police absente ; l'usage privé exige une police adaptée.

Dans le clavier, \key{F2} ou \key{Settings}, puis \key{Input method} :
\begin{choices}
Off — direct keyboard & Saisie américaine ordinaire \\
Romaji & Hiragana japonais depuis une saisie romanisée \\
Pinyin & Chinois à partir du pinyin \\
Zhuyin & Chinois avec la disposition zhuyin de base \\
Wubi 86 & Chinois par codes de forme \\
US-International & Accents latins courants \\
\end{choices}

Choisissez un résultat dans la barre des candidats. Revenez à
\key{Off — direct keyboard} dans le même menu pour la saisie américaine normale.
Le choix s'applique partout et s'enregistre à la fermeture d'Inkline.

\topic{Toucher, stylet et presse-papiers}
\key{Deux doigts verticalement :} parcourez l'historique. Glissez vers le bas
pour voir les sorties anciennes, vers le haut pour revenir vers l'invite.
Chaque terminal garde 500 lignes physiques en RAM, en plus de l'écran courant.

\key{Deux doigts horizontalement :} changez d'un emplacement par geste.
Levez les doigts avant le geste suivant. Écartez-les ou rapprochez-les pour
modifier la taille du texte.

\key{Un doigt :} faites glisser pour sélectionner. Alt/Opt droite+C copie ;
Alt/Opt droite+V colle. Tous les terminaux partagent un presse-papiers en RAM
de 128 KiB maximum.

\key{Stylet :} une application ayant activé les rapports de souris reçoit
les pressions, mouvements et relâchements comme des événements du bouton gauche,
avec le protocole négocié, notamment SGR. Sinon, le glissement sélectionne.
Maintenez \key{Shift} pour imposer une sélection locale dans une application
qui utilise la souris.

\key{OSC 52} permet aux programmes, y compris ceux lancés par SSH, de lire ou
remplacer ce presse-papiers. Il est distinct de celui des carnets.
Alt/Opt droite+X copie et indique que la suppression doit se faire dans l'éditeur :
OSC 52 ne sait pas supprimer son texte.

\key{Images :} les transferts kitty inline et par mémoire partagée restent en
RAM. Les modes fichier et fichier temporaire sont désactivés. Sixel fonctionne
sur demande explicite, mais n'est pas annoncé automatiquement. Les placeholders
kitty, les animations et certains modes sixel restent inachevés.

\code{clear} et l'effacement complet retirent texte et images visibles.
Les images hors écran sont libérées sous pression mémoire. Aucun cache d'images
sur disque ; les grandes images consomment tout de même la RAM limitée.

\topic{Clavier USB et machine à écrire}
Ouvrez \key{Settings → USB (page 3)}. \key{Network} rétablit Ethernet USB ;
\key{Send keyboard} transforme la tablette en clavier pour un ordinateur ;
\key{Receive keyboard} reçoit un clavier filaire par un adaptateur OTG alimenté.
Le Type Folio reste disponible. Installez l'utilitaire Python du site pour
l'envoi au clavier ; aucun module Python supplémentaire n'est nécessaire.

Choisissez Send keyboard, branchez l'USB-C, réglez le profil ci-dessous puis
ouvrez \key{Open typewriter}. Placez le curseur dans le document sur l'ordinateur
et touchez \key{Start}. L'ouverture se fait toujours en pause.
\key{Escape} ou \key{Stop} arrête l'envoi et relâche les touches. Quitter cet écran,
changer de terminal, perdre le focus, ouvrir une application native ou mettre
en veille suspend aussi l'envoi. Aucune reprise automatique après reconnexion.

\begin{choices}
US / ASCII & Disposition US, Verr. Maj désactivé. Texte non ASCII refusé. \\
Mac & Source de saisie Unicode Hex Input, Verr. Maj désactivé. Jusqu'à U+FFFF ;
les caractères des plans supplémentaires sont refusés. \\
Linux & Disposition US, Verr. Maj désactivé. L'application doit accepter
Ctrl+Shift+U, code hexadécimal, Entrée. Toutes ne le permettent pas. \\
Windows & Installez WinCompose sur le PC, Compose sur Right Alt, saisie Unicode
activée. Disposition US, Verr. Maj désactivé. Séquence : Compose, u, six chiffres hexadécimaux, Entrée. \\
\end{choices}
WinCompose : \url{https://wincompose.info/}. Les règles personnalisées ne doivent
pas remplacer cette séquence. HID n'offre pas de canal de texte Unicode universel :
la configuration et les polices du destinataire comptent. Essayez un court texte.

\key{Input method} reprend Off/US direct, Romaji, Pinyin, Zhuyin, Wubi 86 et
US-International. Seuls les caractères validés sont envoyés. Changer de méthode
met en pause ; reprenez avec Start. \key{Off} rétablit la saisie US ordinaire.
\key{Unicode} ou Alt+Space ouvre le sélecteur de caractères pour la sortie USB.

L'aperçu récent reste en RAM ; ce n'est ni un document enregistré ni une copie
de l'éditeur distant. Les flèches et raccourcis agissent sur cet éditeur.
Alt/Opt droit+C copie l'aperçu, V envoie le presse-papiers Inkline. Coupez dans
l'éditeur destinataire. Un accusé USB ne confirme pas le contenu affiché.

\topic{Envoyer un fichier, rétablir le réseau}
Les commandes sont dans \path{~/.local/bin}. En SSH standard, utilisez aussi
\code{~/inkline usb} et \code{~/inkline type}. Les outils Outpost restent intacts.
\UsbExamples
Les fichiers UTF-8 n'ont \key{aucune taille maximale fixe}. La lecture se fait
par petits blocs, sans copie complète en RAM ni fichier temporaire. Les fichiers
repositionnables sont vérifiés en entier avant l'envoi puis relus ; ne les modifiez
pas pendant l'opération. Un pipe est validé au fil de l'eau : une erreur peut
survenir après un envoi partiel. Le BOM UTF-8 initial est retiré ; CRLF/CR deviennent
des retours à la ligne. Entrée et Tab agissent sur l'application au premier plan.

\code{--cps} (20 par défaut) ou \code{--delay-ms} règle la vitesse ; Unicode
nécessite davantage de rapports. \code{--start-delay} laisse le temps de placer
le focus. Dry-run affiche les rapports sans USB. Ctrl+C ou \code{--stop} annule
et relâche les touches. Un seul émetteur à la fois. Une file interactive pleine
suspend l'envoi. Vérifiez le document avant de reprendre un envoi interrompu.
Texte et état temporaire ne sont pas écrits en flash. Le profil est enregistré
avec les autres réglages à la fermeture normale.

\key{Network} ou \code{inkline-usb network} rétablit Ethernet. Quitter Inkline,
même après un plantage ou une récupération d'urgence, rétablit aussi le réseau.
Le mode USB n'est pas mémorisé au démarrage.
\UsbRecoveryExample
Ce minuteur est indépendant du shell. Changez de mode sur la tablette ou en SSH
Wi-Fi, jamais en SSH USB. Un échec tente de rétablir Ethernet. Une session PPP
satellite active ou un gadget inconnu garde le port, y compris lors du nettoyage
à la fermeture. Opt+RightAlt+Backspace reste disponible pour redémarrer Inkline,
mais ferme tous ses terminaux et perd leur travail non enregistré.

\topic{Lancer des programmes au clavier}
Enregistrez les raccourcis depuis un shell Inkline. Les arguments sont littéraux ;
appelez explicitement un shell pour sa syntaxe. Les programmes de raccourci
lisent \code{~/.bashrc} par l'intermédiaire de Bash.

Settings → Command shortcuts (page 2) permet d’associer une commande à une lettre. Choisissez Terminal ou Native e-paper app, puis Save. Remove demande confirmation ; Reset draft annule les modifications. T est verrouillé. Seuls Save et la suppression confirmée écrivent sur le stockage. L’outil en ligne de commande gère les mêmes raccourcis.

Tous les raccourcis globaux utilisent désormais Opt+RightAlt, y compris T et la récupération par maintien de Backspace. Utilisez la touche Opt distincte du Folio ; sur un clavier USB, utilisez Windows/Command (Super). Ctrl+Alt reste disponible dans Emacs. Quittez puis rouvrez Inkline après la mise à jour ; l’installation préserve les terminaux ouverts.

\begin{shell}
~/inkline shortcut register e emacs
~/inkline shortcut register p python3
~/inkline shortcut list
~/inkline shortcut deregister p
\end{shell}

\key{Opt+RightAlt+E} ouvre Emacs dans un emplacement libre.
\code{~/inkline run PROGRAM ARG...} lance sans enregistrer de raccourci.
Un nouvel enregistrement remplace l'ancien ; le retrait ne touche pas aux
programmes en cours. Les touches suivent les positions physiques américaines.
Mettez la ponctuation entre guillemets et utilisez des chemins complets pour
vos programmes, qui doivent déjà être installés.

Pour une application Qt/e-paper native qui dessine son propre écran :
\begin{shell}
~/inkline shortcut register a --epaper /home/root/my-app
\end{shell}
Une seule application native s'exécute à la fois et suspend Inkline et ses shells.
Le lanceur fournit les réglages Qt et des caches en RAM, ignore les journaux et
reprend Inkline à la fin. Utilisez les entrées Qt sans capture exclusive du clavier.

\begin{notice}
\key{Opt+RightAlt+T est permanent.} Impossible de l'enregistrer ou de le retirer.
Il arrête l'application native et retrouve les terminaux existants.
Ne maintenez pas Ctrl pour ce raccourci ; le remappage de Caps Lock reste local à Inkline.
\end{notice}

\key{Urgence :} maintenez Opt+RightAlt+Backspace deux secondes pour arrêter
l'application et ses descendants, puis redémarrer Inkline. Les terminaux
ferment et le travail non enregistré est perdu. Une application root peut
bloquer ce secours en capturant les entrées, modifiant les services ou épuisant
les ressources. Le redémarrage de la tablette reste le dernier recours.

\topic{Installer, mettre à jour, garder le guide}
Le paquet vise \key{reMarkable 2 avec le firmware 3.27} et ses bibliothèques
Qt 6.8 e-paper d'origine. Téléchargements et contrôles d'intégrité :

\InstallLink

L'installeur vérifie matériel, firmware, polices, démarrage de Qt, répertoires
temporaires en RAM et swap désactivé. Les versions sont conservées dans
\path{/home/root/.local/share/inkline} ; le lanceur démarre avec la tablette.
La mise à jour garde les terminaux ouverts. \key{Quittez et rouvrez Inkline
pour utiliser la nouvelle version.} Les anciennes versions occupent de la place
et permettent un retour en arrière.

Ce PDF unique contient les neuf langues et accompagne le paquet. L'installation
tente son importation dans \key{My files} par l'interface Web USB lorsque
celle-ci est activée et que les carnets fonctionnent. Une importation réussie
est mémorisée pour éviter de répéter celle du même guide.

Sinon, branchez l'USB, quittez Inkline et activez \key{USB web interface} dans
les réglages de stockage des carnets. Depuis SSH :
\begin{shell}
~/inkline manual
\end{shell}
Vous pouvez aussi importer le PDF téléchargé par le navigateur USB ou l'application
reMarkable pour ordinateur/téléphone. L'importateur ne modifie jamais directement
la base des carnets et n'interrompt pas les terminaux. Le guide importé reste dans
My files après la désinstallation.

Depuis SSH, revenez aux carnets ou désinstallez Inkline :
\begin{shell}
~/inkline stop
~/inkline uninstall
\end{shell}
L'installation des logiciels facultatifs, dont le profil LaTeX recommandé,
est expliquée sur le site Web.

\topic{Fonctionnement, stockage et RAM}
Qt transmet les touches à un mappage commun des raccourcis et de Caps Lock.
libghostty-vt encode les touches du terminal ; chaque emplacement relie son
programme à un pseudo-terminal (PTY). Les shells reçoivent
\code{HOME=/home/root} ; Bash interactif lit \code{~/.bashrc}.

La sortie traverse l'analyseur du terminal et le décodeur sixel. libghostty
gère cellules, curseur, écrans alternatifs, historique et images kitty.
Le moteur conserve une image grise et redessine les lignes modifiées ; le
backend e-paper Qt met à jour la dalle. Un démon evdev indépendant traite
les raccourcis globaux et l'alimentation ; systemd gère veille et réveil.

Chaque terminal limite les allocations du cœur à 64 MiB, les images kitty à
16 MiB par écran et les images sixel conservées à 16 MiB. Le rendu et les
applications s'ajoutent à ces valeurs. Les emplacements vides n'allouent pas de
terminal. Neuf grosses applications peuvent dépasser la RAM disponible.

Images temporaires, caches et sessions utilisent la RAM. Logiciels installés,
documents, réglages enregistrés à la sortie et fichiers des applications
utilisent le stockage persistant. Inkline n'interdit pas les écritures des programmes.

Le paquet Python facultatif utilise Python standard. Pour réduire les caches,
ajoutez à \code{~/.bashrc} sur la tablette :
\begin{shell}
export PYTHONDONTWRITEBYTECODE=1
export PIP_NO_CACHE_DIR=1
\end{shell}
Exécutez \code{source ~/.bashrc} ou ouvrez un shell. Utilisez
\code{python3 -m pip install --no-compile PACKAGE} pour éviter la compilation
de bytecode propre à l'installation par pip. Le mode isolé de Python ignore
ses variables d'environnement. Le guide Python complet est sur le site.

Inkline est sous licence GPL-3.0-or-later. Les paquets contiennent les sources
correspondantes et les mentions des dépendances. Les comptes rendus séparent
tests automatisés et comportements physiques restant à confirmer.
